% !TeX program = lualatex
\documentclass[theme=modern,subject=Mathematik,grade=8. Klasse,number=2]{schoolexam}

\examsetup{
  title={Klassenarbeit},
  class={8. Klasse},
  duration={90 Minuten},
  aids={Geodreieck, Lineal, Taschenrechner},
  totalpoints={60 Punkte},
  teacher={}
}

\begin{document}

\begin{examtask}[title={Terme und Gleichungen},points=6]
Vereinfache die Terme.\par\medskip
a) \quad $3(x+2)-2(x-5)$ \hfill \answerline[7cm]\par\medskip
b) \quad $4(2x-1)+3(x+4)$ \hfill \answerline[7cm]\par\medskip
c) \quad $5(x-2)-(2x+7)$ \hfill \answerline[7cm]
\end{examtask}

\begin{examtask}[title={Gleichungen lösen},points=8]
Löse die folgenden Gleichungen nach $x$ auf.\par\medskip
\begin{tabularx}{\linewidth}{@{}X X@{}}
a) $2x+5=17$\par $x=$ \answerline[5cm] &
b) $3x-4=11$\par $x=$ \answerline[5cm]\\[7mm]
c) $5x+3=2x+15$\par $x=$ \answerline[5cm] &
d) $4(x-2)=2(x+3)$\par $x=$ \answerline[5cm]
\end{tabularx}
\end{examtask}

\begin{examtask}[title={Sachaufgabe},points=4]
Lea kauft 3 Hefte und 2 Stifte und bezahlt dafür 8,20~€. Ein Heft kostet 1,80~€.
Wie viel kostet ein Stift?\par\medskip
Rechnung:\par\smallskip
\gridbox[28]{3}\par\medskip
Antwort: \answerline[8cm]
\end{examtask}

\examnewpage

\begin{examtask}[title={Geometrie},points=6]
Berechne den Flächeninhalt der folgenden Figuren.\par\medskip
\begin{tabularx}{\linewidth}{@{}X X@{}}
a) Rechteck mit $6\,\mathrm{cm}$ und $4\,\mathrm{cm}$\par\medskip
$A=$ \answerline[5cm]
&
b) Dreieck mit $g=8\,\mathrm{cm}$ und $h=5\,\mathrm{cm}$\par\medskip
$A=$ \answerline[5cm]
\end{tabularx}
\end{examtask}

\begin{examtask}[title={Prozentrechnung},points=4]
In einer Klasse sind 20 Schülerinnen und Schüler. 15~\% davon haben eine Brille.
Wie viele Schülerinnen und Schüler haben keine Brille?\par\medskip
Rechnung:\par\smallskip
\gridbox[24]{3}\par\medskip
Antwort: \answerline[8cm]
\end{examtask}

\begin{examtask}[title={Zuordnungen und Funktionen},points=4]
Gegeben ist die Zuordnung: $y=2x-1$.\par\medskip
a) Berechne $y$ für $x=3$. \hfill \answerline[7cm]\par\medskip
b) Berechne $x$ für $y=7$. \hfill \answerline[7cm]
\end{examtask}

\examnewpage

\begin{examtask}[title={Textaufgabe},points=6]
Ein Fahrrad kostet ursprünglich 480~€. Es wird zuerst um 15~\% reduziert und nach
einigen Wochen nochmals um 10~\%. Wie viel kostet das Fahrrad am Ende?\par\medskip
Rechnung:\par\smallskip
\gridbox[28]{3}\par\medskip
Antwort: \answerline[8cm]
\end{examtask}

\begin{examtask}[title={Geometrie -- Winkel},points=4]
In einem Dreieck beträgt ein Innenwinkel $45^\circ$ und ein anderer $70^\circ$.
Berechne die Größe des dritten Innenwinkels.\par\medskip
Rechnung:\par\smallskip
\gridbox[28]{3}\par\medskip
Antwort: \answerline[8cm]
\end{examtask}

\vfill
\begin{schoolSoftBox}[colback=schoolSoft,colframe=schoolSoft,arc=2.5mm]{\faTrophy\hspace{2mm}Geschafft!}
Überprüfe deine Lösungen noch einmal.\par\smallskip
Erreichte Punktzahl: \answerline[18mm] / 60\hfill Note: \answerline[24mm]
\end{schoolSoftBox}

\end{document}
